% `tex.web` §846: the line breaker's own narration.
%
%     @firstpass
%     []\rip AAA AAA AAA AAA
%     @ via @@0 b=0 p=0 d=0
%     @@1: line 1.2 t=0 -> @@0
%     AAA AAA AAA AAA
%     @\par via @@1 b=0 p=-10000 d=0
%     @@2: line 2.2- t=0 -> @@1
%
% Four parts. §863/§866/§873 head each pass. §857 shows the material since the
% last reported breakpoint in §174's short form — characters bare, a font
% announced by its identifier the first time it is used, `[]` for a box, `|` for
% a rule, a space for any non-zero glue. §855 reports each feasible break: the
% node broken at, the active node reached through, the badness (`*` past
% `inf_bad`), the penalty, and the demerits of *this line* before the running
% total is added. §856 reports each active node created: its serial, the line it
% ends and its fitness class, `-` if the break is at a discretionary, the total
% demerits, and the serial it came from.
%
% The serials are §845's, numbered as passive nodes are created, and they are
% what made this the last of the seven: TeX numbers them from its active list,
% and our breaker keeps a DP table. The table does carry a §835-pruned active
% list, so the numbering is reconstructible — but TeX walks that list oldest
% first where we walk it newest first, so the `@ via` lines are collected and
% reversed before they are printed.
%
% **What this case does not cover, and why.** A paragraph that gets past the
% first pass narrates differently, though it *sets the same*:
%
%     tex:   @firstpass @secondpass                   breaks with `b=*`, `d=*`
%     ours:  @firstpass @secondpass @emergencypass    different `@` lines
%
% TeX has two passes and the second *is* the final one — §854's
% `artificial_demerits` lets it accept an overfull line at zero cost, which is
% what `d=*` means. We have three: our second pass refuses overfull lines and a
% separate emergency pass allows them, so it reaches the same breaks by a
% different route and reports a different search.
%
% The breaks themselves agree. For
%
%     \hsize=90pt \tolerance=9999
%     AA BBBB CC DDDDDD EE FFFF GG HHHHHH II JJ KKKK LL\par
%
% where no break is feasible at any tolerance, both engines set
% `AABBBBCCDDDDDD | EEFFFFGGHHHHHH | IIJJKKKKLL`, every glyph at the same
% horizontal *and* vertical position. So this is a trace difference, not a
% typesetting one — which is why the case stays on a paragraph that breaks in
% the first pass rather than claiming a fault that is not there.
\catcode`\{=1 \catcode`\}=2 \catcode`\#=6
\font\rip=cmr10 \rip
\hsize=100pt \parindent=0pt \tolerance=10000 \hbadness=10000
\pretolerance=100 \baselineskip=12pt
\tracingparagraphs=1
AAA AAA AAA AAA AAA AAA AAA AAA\par
% a second paragraph, so the serials restart per paragraph as TeX's do
AAA AAA AAA AAA\par
\tracingparagraphs=0
\end
